\chapter{BACKGROUND AND RELATED WORK}

This chapter establishes the conceptual foundation of the project. It first frames the bus arrival time prediction problem and the main families of solution approaches, then introduces the GTFS standard and real-time transit data, summarizes the machine learning and deep learning models used in this work, and finally discusses the reference study that motivated the project.

\section{Bus Arrival Time Prediction}

Bus Arrival Time Prediction (BATP) estimates when a bus will reach a given stop, based on its current state and the conditions along its route. The quality of a prediction is governed by several interacting factors: the time of day and day of the week (which capture demand and congestion patterns), the spatial position along the route, dwell time at stops for passenger boarding and alighting, traffic conditions, and weather.

Solution approaches are commonly divided into two families. \textbf{Model-based (analytical)} methods rely on explicit assumptions about traffic flow or use techniques such as Kalman filtering to propagate the state of a bus forward in time. These methods are interpretable and require little training data, but they struggle to capture the highly nonlinear and context-dependent nature of urban travel times. \textbf{Data-driven} methods instead learn the relationship between observable features and travel time directly from historical observations. With the increasing availability of automatic vehicle location data, data-driven machine learning approaches have become the dominant paradigm because they can exploit large numbers of heterogeneous features without hand-crafted traffic models.

A recurring theme in the BATP literature is that travel time is best predicted at the level of individual stop-to-stop \emph{segments}, which are then accumulated to obtain an arrival-time estimate at any downstream stop. This project follows the segment-level formulation, since it isolates the local factors that affect each portion of the route and produces a flexible building block for arrival estimation.

\section{GTFS and Real-Time Transit Data}

The General Transit Feed Specification (GTFS) \cite{gtfs} is the de facto open standard for describing public transportation schedules and the associated geographic information. A \emph{static} GTFS feed is a collection of text files describing agencies, routes, trips, stops, stop sequences, and scheduled arrival and departure times. A complementary specification, GTFS-Realtime, describes live vehicle positions, trip updates, and service alerts.

Static GTFS data is central to this project. The scheduled travel time between two consecutive stops, derived from the \texttt{stop\_times} file, provides a strong prior for how long a segment \emph{should} take. Using this scheduled time, and the deviation of the observed time from it, as explicit model inputs is a key design hypothesis of this work.

The city of İzmir publishes static GTFS data and exposes a public API for live bus positions, but it does not provide a GTFS-Realtime feed. The live API returns the current coordinates of active buses, which must be matched to the nearest stop and accumulated over time to reconstruct trips. This data characteristic, together with the coarse polling interval, directly shapes the methodology and the achievable accuracy, as discussed in later chapters.

\section{Machine Learning Models for Travel-Time Prediction}

Three families of supervised models are employed and compared in this project, ranging from interpretable ensembles to sequence-aware neural networks.

\textbf{Random Forest.} The Random Forest \cite{breiman2001} is an ensemble of decision trees trained on bootstrap samples of the data, with predictions averaged across the trees. It handles nonlinear interactions and mixed feature types robustly, requires little tuning, and provides feature-importance estimates, making it a strong and interpretable baseline for tabular travel-time data.

\textbf{Gradient-Boosted Trees (XGBoost).} XGBoost \cite{chen2016} builds an additive ensemble of trees in which each new tree corrects the residual errors of the previous ones. With regularization and efficient training, gradient boosting frequently achieves state-of-the-art accuracy on tabular problems. In this project XGBoost serves as the strongest classical model.

\textbf{Long Short-Term Memory Networks.} For sequential data, recurrent neural networks model temporal dependencies, and the Long Short-Term Memory (LSTM) architecture \cite{hochreiter1997} mitigates the vanishing-gradient problem through gated memory cells. Because consecutive segments within a trip form a natural sequence, an LSTM can in principle capture how recent travel-time deviations propagate forward. This project uses a dual-input LSTM that processes a window of recent segments together with static contextual features.

For all of these models, the predictive performance depends heavily on \emph{feature engineering}: the construction of temporal, spatial, schedule-based, historical, and dwell-time features that expose the relevant signal to the learner. The design and contribution of these features is a central part of the methodology in Chapter~3.

\section{The Reference Study}

The work of Kaya and Kalay \cite{kaya2025} is the direct reference for this project. They proposed context-aware deep learning models for the spatio-temporal forecasting of bus arrival times in urban transit systems, incorporating contextual information into the prediction and reporting a Mean Absolute Error of 2.97 minutes.

This project shares the central idea of context-aware prediction but differs from the reference study in three important respects. First, it is built on an \emph{independently collected} real-time dataset from the İzmir ESHOT network, rather than a pre-existing dataset. Second, it explicitly compares classical machine learning models against a deep learning model on the \emph{same} test set and quantifies the difference with statistical significance tests, instead of focusing on a single model family. Third, it adopts an A/B-driven methodology in which every proposed enhancement must demonstrate a measurable improvement to be retained; as Chapter~5 shows, this rigor reveals that several commonly assumed improvements do not in fact help, and that the achievable error is bounded by a data-measurement floor.
